\section{Layer L2: formal semantics with RDF, OWL, SHACL and SPARQL}
\label{sec:layer2}

L1 tells an agent what things \emph{are}. L2 tells it what \emph{follows}, what
is \emph{forbidden}, and how to \emph{ask}. This is the layer
industrial practice most often skips, and the experiments in
\S\ref{sec:results} suggest that skipping it caps agent accuracy at roughly the
\CeilingCTwo{} mark regardless of how good the type model is.

\subsection{RDF and RDFS: the join substrate}

RDF \cite{w3c2014rdf} represents everything as subject--predicate--object
triples over IRIs and literals. Its virtues for agents are unglamorous and large.

\emph{Schemaless join.} Merging an \opcua{} address space, an AAS environment, a
maintenance log and a recipe database requires no shared schema and no ETL
target, only shared IRIs. Our reference graph does exactly this: signals from
the control system, nameplates from the AAS, topology from the engineering model,
events from the historian and recipes from the MES, in \NTriples{} triples with
no integration schema.

\emph{Open-world assumption.} An agent working over a plant model must
distinguish ``no interlock is recorded'' from ``there is no interlock''. RDF's
open world makes the distinction expressible; closed-world SQL thinking makes it
disappear, and the disappearance is dangerous in exactly the safety-relevant
cases.

\emph{Named graphs and quads} give per-source provenance for free, which is what
makes an agent's evidence auditable.

\paragraph{Limitation.} RDF alone says nothing about units, behaviour or
constraints (U3, U5, U6 all $0$), and its verbose serialisations are the
worst-in-class on U9: expanded JSON-LD and N-Triples measured at
$64\times$ and $\MaxSerialRatio\times$ the tag-list token cost respectively
(Table~\ref{tab:e3-serialisation}). Turtle is the only serialisation we would
put anywhere near a context window, and even then only in projected slices.

\subsection{OWL 2: inference and its limits}

OWL~2 \cite{w3c2012owl} adds class and property axioms with a
description-logic semantics: subsumption, equivalence, disjointness, domains and
ranges, cardinality, property characteristics (transitive, symmetric, inverse,
functional). Its
profiles trade expressivity for tractability: EL for large terminologies, QL
for query rewriting over databases, RL for rule-based forward chaining.

\paragraph{Agent affordances.} Two matter. \emph{Materialised inference}
converts implicit facts into retrievable ones: declaring
\texttt{isPartOf} transitive means ``every asset in the site'' is a query rather
than a recursive traversal the agent must implement. Declaring
\texttt{feeds}~/~\texttt{isFedBy} inverse means impact analysis works in both
directions from one asserted set. \emph{Consistency checking} catches
contradictory assertions, so an agent that concludes a node is both a
\texttt{Sensor} and an \texttt{Actuator} can be told it is wrong, not merely
scored as unlikely.

\paragraph{Where OWL is the wrong tool.} OWL's open-world, no-unique-name
semantics makes it a poor \emph{validator}. ``This signal must have exactly one
unit'' is, in OWL, not a constraint that can be violated but an axiom from which
the reasoner will happily infer that two differently-named units are the same
thing. Practitioners have repeatedly rediscovered this; it is the reason SHACL
exists, and it is why we use OWL for the T-Box and SHACL for validation in our
implementation.

\subsection{SHACL: turning generation into a checkable proposition}
\label{sec:shacl}

SHACL \cite{knublauch2017shacl} (W3C Recommendation, 2017) validates RDF graphs
against \emph{shapes}:
node shapes with target declarations and property shapes carrying constraint
components (\texttt{minCount}, \texttt{datatype}, \texttt{in},
\texttt{minInclusive}, \texttt{class}, \texttt{node}, \texttt{closed}), plus
\texttt{sh:sparql} for constraints that need the full query language. Validation
returns a report: focus node, path, severity, message.

This is, in our assessment, the most under-exploited technology in the entire
survey for agent work, and the experimental centrepiece of this paper. The
pattern is simple and general:

\begin{enumerate}[label=\textbf{\arabic*.}]
  \item the agent proposes an action or an answer;
  \item the proposal is \emph{lifted} into RDF against the plant graph;
  \item SHACL validates the union;
  \item a conforming proposal executes; a violating one is returned to the model
  \emph{with the validation report as a repair signal}.
\end{enumerate}

Generation stays statistical; acceptance becomes deductive. The report is not
merely a rejection: \texttt{sh:message} with interpolated result variables
produces a natural-language explanation the model can act on, which is why we
report \emph{diagnosis rate} alongside detection rate in Experiment~E4. A
guardrail that says only ``no'' is far less useful in a repair loop than one
that says \emph{which} constraint failed and by how much.

Listing~\ref{lst:shacl} shows the constraint that catches the most dangerous
fault family in our corpus. Note that it is not expressible without \emph{both}
L2 (SPARQL over a graph) and L3 (QUDT conversion factors), a concrete instance
of the non-substitutability claim of \S\ref{sec:framework}.

\begin{lstlisting}[language=turtle,float=*,floatplacement=t,caption={A unit-aware
range constraint. The proposed value is converted into the signal's own unit
through the QUDT affine conversion before the engineering range is applied. A
validator without unit semantics either accepts a value that is out of range (a
false negative) or rejects a legal value expressed in a different unit (a false
alarm). In Experiment~E4, tier~G2 does both.},label={lst:shacl}]
shp:SetpointWriteShape a sh:NodeShape ;
  sh:targetClass sio:SetpointWrite ;
  sh:sparql [ a sh:SPARQLConstraint ;
    sh:message "F-RANGE: {?conv} is outside [{?lo},{?hi}] of {?sigtag}" ;
    sh:select """
      SELECT $this ?conv ?lo ?hi ?sigtag WHERE {
        $this sio:targetSignal ?sig ; sio:value ?v ; sio:valueUnit ?u .
        ?sig qudt:unit ?su ; sio:engineeringLow ?lo ;
             sio:engineeringHigh ?hi ; skos:notation ?sigtag .
        ?u  qudt:conversionMultiplier ?mu ; qudt:conversionOffset ?ou ;
            sio:dimensionVector ?du .
        ?su qudt:conversionMultiplier ?ms ; qudt:conversionOffset ?os ;
            sio:dimensionVector ?ds .
        FILTER (?du = ?ds)
        BIND ((((?v * ?mu) + ?ou) - ?os) / ?ms AS ?conv)
        FILTER (?conv < ?lo || ?conv > ?hi) } """ ] .
\end{lstlisting}

\subsection{SPARQL: the aggregation escape hatch}

SPARQL~1.1 \cite{w3c2013sparql} supplies pattern matching, \texttt{OPTIONAL},
negation, property paths (including transitive \texttt{+} and \texttt{*}),
aggregation, subqueries, federation and \texttt{ASK}. For an agent it plays three distinct roles.

\emph{Retrieval.} A SPARQL \texttt{CONSTRUCT} or \texttt{DESCRIBE} produces a
faithful, minimal subgraph. This is the projection operation that L1 formats lack
and that U9 demands.

\emph{Computation.} Aggregations, joins and transitive closures are evaluated by
the engine, not by the model. Experiment~E7 quantifies this: eight questions that
consumed $5{,}000$--$5{,}500$ tokens of stuffed context are answered by a query
plus its result set in $79$--$314$ tokens, a saving of up to
$\MaxQuerySaving\times$. One of them, a plant-wide alarm sweep, is not
answerable by stuffing at \emph{any} budget while costing about two hundred
tokens as a query.

\emph{Verification.} An \texttt{ASK} query is a truth test. An agent that says
``the compressor feeds the Line~1 filler'' can have the claim checked in one
round trip.

\paragraph{The text-to-SPARQL problem.} All of this presumes the agent can write
the query. Query generation over an unfamiliar schema is itself an accuracy
problem, well studied for SQL and less so for SPARQL, and it is the main risk of
the query-based architecture. The mitigation we recommend, and implement in the
artefact, is not free-form generation but \emph{parameterised query templates
exposed as tools}: the agent chooses a template and binds arguments, so the
argument space is constrained (Experiment~E5) and the query is correct by
construction.

\subsection{PROV-O: provenance for agent writes}

PROV-O \cite{w3c2013prov} supplies the vocabulary for U7: activities, agents,
entities, \texttt{wasDerivedFrom}, \texttt{wasAssociatedWith},
\texttt{startedAtTime}. Its importance rises sharply the moment agents write as
well as read: an autonomous change to a plant needs an audit trail with the same
standing as an operator's signature.

It is the only one of the further L2 vocabularies we score separately
(Table~\ref{tab:auf-matrix}); the rest are enabling tooling rather than distinct
agent-usability entries. Three matter in a deployment: \emph{SKOS}
(\texttt{skos:altLabel}/\texttt{skos:notation} are the synonymy fields our entity
resolver reads), \emph{JSON-LD} (the JSON/RDF bridge that lets WoT~TD and NGSI-LD
be developer-friendly and grounded at once), and the \emph{R2RML/OBDA} mapping
languages that keep the graph \emph{virtual} over the historian --- the realistic
deployment path when plant data is too large to materialise as triples.

\subsection{Summary of L2}

L2 contributes precisely the capabilities that L1 lacks and that our experiments
show to be decisive: transitive relational reasoning (U4), checkable constraints
(U6), provenance (U7) and aggregation pushdown (U8). It contributes nothing on
units (U3), which is L3's job, and it is expensive on U9 unless used as a
query interface rather than a serialisation. The single most actionable
recommendation of this survey is that the L2 layer should be adopted for
\emph{validation and query}, even by organisations that have no intention of
replacing their L1 models with RDF.
